% ============================================================
%  3. Distribuirani sustavi i tolerancija na greške
%  DUBINSKO poglavlje · budžet 11 str.
%  Struktura je zaključana.
%  CAP je namjerno UNUTAR 3.5, nema vlastiti odjeljak:
%  ne služi ničemu nizvodno osim 9.1.
% ============================================================
\chapter{Distribuirani sustavi i tolerancija na greške}
\label{ch:distribuirani}

\todo{Uvodni odlomak poglavlja: što poglavlje tvrdi i čemu služi nizvodno.}

% ------------------------------------------------------------
\section{Distribuirani sustav i djelomični otkaz}
\label{sec:djelomicni-otkaz}
% Budžet: 3 str.
% TVRDI: udaljeni poziv se ne smije modelirati kao lokalni, jer uvodi
%        klasu otkaza koja lokalno ne postoji — djelomični otkaz.
% NOSIVI IZVOR: Waldo i dr. 1994 (SMLI TR-94-29) — čitati u cijelosti.
% UZ NJEGA: zablude o distribuiranom računarstvu.
%   ⚠ PROVENIJENCIJA JE SPORNA: sam Deutsch (SE Radio ep. 470, 2021)
%     osporava uvriježenu kronologiju "Deutsch 1994 + Gosling 1997".
%     Napisati pošteno — od zamke na obrani postaje poen za temeljitost.
% NIZVODNO: 6.1 (granica ERP↔posrednik je točno takva granica)

\todo{Definicija distribuiranog sustava; djelomični otkaz kao klasa koja
u lokalnom pozivu ne postoji; zablude o distribuiranom računarstvu uz
poštenu provenijenciju.}

% ------------------------------------------------------------
\section{Pojmovi: uzrok, pogreška i otkaz}
\label{sec:fault-error-failure}
% Budžet: 1,5 str.
% TVRDI: bez razlikovanja uzroka (fault), pogreške (error) i otkaza
%        (failure) rad ne može dosljedno govoriti o "načinima otkazivanja".
% NOSIVI IZVOR: Avizienis i dr. 2004, IEEE TDSC 1(1) — u cijelosti.
% NIZVODNO: daje JEZIK cijelom poglavlju 6.2. Bez ovog odjeljka izvod
%           u 6.2 nema dosljednu terminologiju.

\todo{Taksonomija fault/error/failure; hrvatski termini i njihovo
dosljedno korištenje kroz ostatak rada.}

% ------------------------------------------------------------
\section{Klasifikacija otkaza}
\label{sec:klasifikacija-otkaza}
% Budžet: 2 str.
% TVRDI: otkazi se dijele na crash / omission / timing / Byzantine;
%        fiskalizacijski tijek pogađaju prve tri. Većina softverskih
%        kvarova je PROLAZNA — i to je ono što opravdava retry.
% IZVORI: Cristian 1991, CACM 34(2) — u cijelosti.
%         ⚠ ACM metapodaci bilježe "Flavin Cristian"; ispravno je FLAVIU.
%         Gray 1985/1986, "Why Do Computers Stop" — ciljano (Heisenbugs).
% NIZVODNO: 6.2 (koje klase uzimamo u razmatranje), 4.1 (opravdanje retryja)

\todo{Klase otkaza; koje su relevantne za fiskalizacijski tijek i zašto
bizantski otkaz nije; prolaznost kvarova kao opravdanje ponovnog pokušaja.}

% ------------------------------------------------------------
\section{Timeout ne razlikuje sporo od mrtvog}
\label{sec:detekcija-otkaza}
% Budžet: 1,5 str.
% TVRDI: nepouzdana detekcija otkaza je TEMELJNO ograničenje, ne
%        implementacijski nedostatak. Posljedica koju rad koristi svuda:
%        "nema odgovora" nikad ne znači "nije zaprimljeno".
% IZVOR: Chandra & Toueg 1996, JACM 43(2) — CILJANO.
%        Samo definicije potpunosti i točnosti detektora, ne dokazi.
% NIZVODNO: 6.2 (sinkroni SOAP prema CIS-u nema propisan timeout),
%           4.1 (zašto timeout sam nije rješenje)

\todo{Detektori otkazâ; potpunost i točnost; nemogućnost razlikovanja
sporog od mrtvog procesa i što to znači za tumačenje izostalog odgovora.}

% ------------------------------------------------------------
\section{Semantika dostave i nemogućnost exactly-once}
\label{sec:semantika-dostave}
% Budžet: 3 str. (2 str. + 1 str. za CAP, koji je ovdje apsorbiran)
%
% ★ NAJVRJEDNIJI ODJELJAK POGLAVLJA.
% TVRDI: exactly-once DOSTAVA ne postoji; ono što se u praksi zove
%        exactly-once je at-least-once + idempotentnost, tj. exactly-once
%        UČINAK. Sam lanac citiranja je doprinos ovog poglavlja.
%
% ⚠ NE POSTOJI rad koji dokazuje "exactly-once dostava je nemoguća".
%   Tko takav citira, citira blog. Pošteni lanac:
%     Gray 1978, §5.8.3.3.1 "THE GENERALS PARADOX"  (ciljano — provjereno u PDF-u)
%       → Saltzer, Reed, Clark 1984, end-to-end argument (u cijelosti)
%          → duplikati se MORAJU rješavati na krajevima, ne u transportu
%       → Helland 2012, "Idempotence Is Not a Medical Condition" (u cijelosti)
%          citirati ACM Queue verziju kao izvornu
%       → Wang i dr. 2021, SIGMOD (ciljano) — kako se to zapravo radi
%   FLP (Fischer, Lynch, Paterson 1985) — SAMO CITIRATI, ne čitati u cijelosti.
%   Halpern & Moses 1990 — IZBAČENO: Gray 1978 daje isti
%     argument čitljivije i njega smo stvarno pročitali.
%
% CAP, APSORBIRAN OVDJE KAO ODLOMAK (ne odjeljak):
%   Gilbert & Lynch 2002 (u cijelosti) — formalni dokaz.
%   Abadi 2012, PACELC — samo citirati: kompromis postoji i BEZ particije.
%   ⚠ Često citirani MIT link .../Gilbert/Brewer2.pdf JE RAD IZ 2012.
%     Rad iz 2002. je .../Gilbert/Brewer6.pdf.
%
% NIZVODNO: 6.3 (arhitektura), RQ2 ("što ostaje nesvodivo")

\todo{Semantika dostave: at-most-once, at-least-once, exactly-once.
Zašto exactly-once dostava ne postoji i kako se postiže exactly-once
učinak. Idempotentnost. CAP i PACELC u tri rečenice.}
